\documentclass[11pt]{article}
\usepackage[margin=1in]{geometry}
\usepackage{amsmath,amssymb,mathtools}
\usepackage{fontspec,unicode-math}
\setmainfont{Latin Modern Roman}
\setsansfont{Latin Modern Sans}
\setmathfont{Latin Modern Math}
\usepackage{xurl,hyperref}
\hypersetup{colorlinks=true,urlcolor=blue,linkcolor=blue}
\newcommand{\sref}[2]{\hyperlink{source:#1:#2}{[S#2]}}
\newcommand{\cataloguescope}[1]{\paragraph{Recorded mathematical question.}#1}
\begin{document}
% UnsolvedMath ID 10000020; source catalogue code AMR-099-0020
\setcounter{section}{10000019}
\section{Exponential intersection tails for loop-erased random walk}\label{10000020}
{\small\sffamily UnsolvedMath ID 10000020 · Probability}
\subsection{Definitions and mathematical statement}

\paragraph{Shared notation.}$\mathbb N=\{1,2,\ldots\}$, and $\mathbb Z,\mathbb Q,\mathbb R,\mathbb C$ denote the integers, rationals, reals and complex numbers. Any other conventions are specified locally.


For an integer $d\ge1$, let $\mathbb Z^d$ have nearest-neighbor edges and root $0$. Chronological loop erasure of a finite path successively deletes each loop as soon as it is completed. Let $\mu_d$ be the infinite loop-erased random-walk law from $0$, defined as the local weak limit of loop erasures of simple random walks stopped on first leaving successively larger boxes. For two independent paths $\gamma_1,\gamma_2$ with law $\mu_d$, let $|\gamma_1\cap\gamma_2|$ count distinct common edges.
\paragraph{Statement recorded in the supplied source.}
For which dimensions $d$ does $\mu_d$ have exponential intersection tails: is there $0<\theta<1$ such that $\mathbb P_{\mu_d\otimes\mu_d}(|\gamma_1\cap\gamma_2|>n)\le\theta^n$ for every integer $n\ge1$? \sref{10000020}{2}
\subsection{Short English statement}
Determine whether infinite loop-erased random walks have exponentially small probabilities of sharing many edges.
\subsection{Sources}
\begin{description}
\item[\hypertarget{source:10000020:1}{S1}] ULAM.AI and the UnsolvedMath Contributors, \emph{UnsolvedMath: A Curated Collection of Open Mathematics Problems}, record 10000020 (AMR-099-0020). CC BY 4.0. \url{https://huggingface.co/datasets/ulamai/UnsolvedMath}.
\item[\hypertarget{source:10000020:2}{S2}] Itai Benjamini, \emph{Coarse Geometry and Randomness}, Saint-Flour lecture notes, version dated 30 October 2013; original author manuscript preserved in Arquivo.pt capture of 31 December 2020, Definition 4.16 and Open Problem 4.22, pp. 35--36. \url{https://arquivo.pt/noFrame/replay/20201231041548id_/http://www.wisdom.weizmann.ac.il/~itai/stflouraug24.pdf}.
\end{description}
\label{end:10000020}
\end{document}
